
\newpage

\section{Proof by Evaluation}\label{sec:pbe}

Our terms are measure application, constructor application, and variables.
\[e ::= x \mid C\;\bar e \mid m\;\bar e\]
Measures are represented as a list of guarded lambdas $[\left<p,\lambda x.e\right>,\ldots]$,
where $p$ is predicate that guards evaluation and the type checker can use to prove termination

We have the following contexts:
\[\Gamma ::= \bullet\mid \Gamma; i :T \mid \Gamma; m : \left<p,\lambda x.e\right> \qquad
\Delta ::= \bullet\mid \Delta; \phi \qquad
\Sigma ::= \bullet\mid \Sigma; \star \mid \Sigma; t\]
$\Gamma$, has our measures, and carries around refinement
type information for free variables in scope. $\Delta$ just exists to
dump all the predicates into, and $\Sigma$ hold terms awaiting PBE. $\star$ in
$\Sigma$ runs PBE on the term being scrutinized.

We have the following operational semantics:
\[
\inference{\left<\Gamma,e_i\right> \hookrightarrow \left<\Gamma,e_i'\right>}{\left<\Gamma,C \bar e\right> \hookrightarrow\left<\Gamma,C\bar e'\right>}
\qquad
\left<\Gamma,x\right> \hookrightarrow\left<\Gamma,x\right>
\]
\[
\inference{\Gamma \vdash p(x)}
       %----------------------------
{\left<\Gamma;(m:\left<p,\lambda \bar x.e\right>),m\;\bar e\right> \hookrightarrow\left<\Gamma;(m:\left<p,\lambda \bar x . e\right>),e[\bar e/\bar x]\right>}
\]
We proceed to define the judgments for $\vdash$, $\vdash_n$, and $\vdash^*$:

$\vdash p$ is just EUF, so we have the following rules:


\[
\inference        {\vdash a=b}
       %----------------------------
               {\vdash b=a}          [sym]
\qquad
\inference        {}
       %----------------------------
               {\vdash a=a}          [refl]
\]

\[
\inference        {\vdash a=b}
       %----------------------------
            {\vdash f(a)=f(b)}       [cong]
\qquad
\inference  {\vdash a=b \qquad \vdash b = c}
         %-------------------------------
                 {\vdash a=c}       [trans]
\]
Checking equality proofs using these rules is done with a data structure called
an E-dag\cite{nelson-oppen}. We can construct proof terms by abstracting away
the mathematical properties of this data structure.

If we consider proofs that
don't use congruence, then we can model this set of equalities as a graph with
terms as nodes and equalities as edges. This graph is symmetric, has all
identity loops, and transitivity means we can compose paths, so we can think of
paths in this graph as morphisms of a groupoid in which objects are terms and
morphisms are equalities --- each of these three rules correspond to one of the
groupoid laws.

The advantage of the groupoid apporach here is that we now have
a way to explicitly construct and compose proof terms, where we can use any of
the equalities in $G$ to give us a literal proof term, construct morphisms
$id_a: a \rightarrow a$ using $[refl]$, construct inverses $p:a\rightarrow b$
from $q:b\rightarrow a$, using $[sym]$, and compose morphisms using $[trans]$.

Congruence, then, is given by the family of endofunctors
$C_-$, indexed over functions in our logic, and is defined in the following equations
\begin{align*}
  C_f (t) &= f t \\
  C_f (p_{a=b}) &= p_{fa = fb}\\
\end{align*}
where $p_{a=b}:a\rightarrow b$ is the morphism that correponds to the equality $a=b$
Note that $C_f(p)^{-1} = C_f(p^{-1})$ and $C_f(p)\circ C_f(q) = C_f(p\circ q)$, in accordance with the Functor laws.

\(\Gamma, \Delta \vdash_n p\) has semantics:
\[
\inference        {\vdash p}
       %----------------------------
           {\Gamma,\Delta\vdash_0 p}      [[0]]
\qquad
\inference
      {\Gamma,\Delta\vdash_0 \phi \rightarrow p}
     %------------------------------------------
       {\Gamma,\Delta;\phi\vdash_0 p}             [[$\Delta$]]
\]
\[
\inference
        {\vdash \Delta \rightarrow q(\bar e) \qquad \Gamma;m:\left<q,\lambda \bar x.e\right>,\Delta;m\;\bar e = e[\bar e/\bar x]\vdash_n p}
     %------------------------------------------------------------------------------
           {\Gamma;m:\left<q,\lambda \bar x.e\right>,\Delta\vdash_{n+1} p}                  [[n]]
\]
Furthermore, we write $\Gamma,\Delta \vdash_* \phi$ if $\exists n | \Gamma,\Delta \vdash_n \phi$

$\Gamma, \Delta, \Sigma \vdash^* p$ defines our PBE semantics and is more involved:
\[
\inference
        {\Gamma, \Delta, \Sigma;\phi \vdash^* \phi}
     %-----------------------------------------------------------------
        {\Gamma,\Delta,\Sigma;\star \vdash^* \phi}
        \qquad
\inference
      {\Gamma;x:\{x:T|R(v)\},\Delta;R(x),\Sigma \vdash^* \phi}
     %--------------------------------------------------------
      {\Gamma;x:\{x:T|R(v)\},\Delta,\Sigma;x \vdash^* \phi}
\]
\[
\inference
        {\vdash \Delta \rightarrow \phi}
     %-------------------------------------
      {\Gamma,\Delta,\Sigma \vdash^* \phi}
      \qquad
\inference
      {\Gamma;C:\bar x \rightarrow \{v:T|R(\bar x)\},\Delta;R(\bar e),\Sigma;\bar e \vdash^* \phi}
     %---------------------------------------------------------------------------------------
       {\Gamma;C:\bar x \rightarrow \{v:T|R(\bar x)\},\Delta,\Sigma;C\;\bar e \vdash^* \phi}
\]
%Add guard for if statements to make sure we terminate by only instantiating measures that actually matter
%this guard should be a predicate in the measure checked by both PBE as well as =_n
\[
\inference
        {\Gamma;m:\left<q,\lambda \bar x.e\right>,\Delta,\Sigma;\bar e \vdash^* q(\bar e) \qquad \Gamma;m:\left<q,\lambda \bar x.e\right>,\Delta;m\;\bar e = e[\bar e/\bar x], \Sigma;e[\bar e/\bar x]\vdash^* p}
     %------------------------------------------------------------------------------
           {\Gamma;m:\left<q,\lambda \bar x.e\right>,\Delta,\Sigma;m\; \bar e\vdash^* p}
\]
\begin{lemma}
  If $G,a=b\vdash u=v$ and then either $a=b$ does not appear as a premise in the proof, or
  there exists $a',b'$ which are subterms of $G$, $u$, or $v$, such that
  $G\vdash a'=a$ and $G\vdash b'=b$ \end{lemma}
\begin{proof}
  Let $p_{ans}$ be the proof term that witnesses the proof $G,a=b\vdash u=v$.

  We proceed by induction on number of occurrences of $C_-$ in $p_{ans}$. In the
  case where we have 0 occurrences, $p_{ans} = id_u \circ p_0\circ p_1 \circ \cdots
  \circ p_k \circ id_v$, where $p_0$ through $p_k$ are equalities from $G$ (or their inverses).
  If $p_{a=b}$ is precomposed in $p_{ans}$ with a $p_i$, then $a\in G$, otherwise, $a=u$,
  and the same is true of postcomposition and $b$, so we just have $a'\equiv a$
  and $b'\equiv b$.

  Now consider the case in which we have $n+1$ occurrences of $C_-$ in
  $p_{ans}$. Let $p_n:\bar u\rightarrow \bar v$ have at most $n$ occurrences of $C_-$.
  By the induction hypothesis, each of $a',b'$ are either subterms of $G$ (in
  which case we are done) or of $\bar u$ or $\bar v$. This problem is
  symmetric, we we really just need to show that this for $a'$ in the
  precomposition case, as before. We have three cases:
  \begin{itemize}
    \item $p_{ans} = id_u\circ C_f(p_n)\circ \cdots$ in which case $u=\bar u$ and we are done
    \item $p_{ans} = \cdots \circ p \circ C_f(p_n)\circ \cdots$ where $p$ or
      $p^{-1}$ are from $G$. Here, $p : \bar u \rightarrow t$, so
      $\bar u$ occurs in $G$.
    \item $p_{ans} = \cdots \circ C_f(p_{n_1}) \circ C_f(p_{n_2} )\circ \cdots$
      where instead of $p_n$, we have $p_{n_i}$ which have at most $n_i$
      occurrences of $C_-$, such that $n_1+n_2=1$. From the functor laws, we
      have $p_{ans} = \cdots \circ C_f(p_{n_1}\circ p_{n_2}) \circ \cdots$,
      which reduces to one of the previous cases, after repeating this case at
      most finitely many times.
\end{itemize}
\end{proof}
\begin{corollary}
  If $G,a=b\vdash u=v$ and $G\not\vdash u=v$, then there exists $a',b'$ which
  are subterms of $G$, $u$, or $v$, such that $G\vdash a'=a$ and $G\vdash b'=b$
  \label{cor:subterm-exist}
\end{corollary}
\begin{proof}
  Let $p_{ans}$ be the proof term that witnesses the proof $G,a=b\vdash u=v$.
  $p_{a=b}$ must appear in $p_{ans}$, since otherwise $G\vdash u=v$.
\end{proof}

\begin{lemma}
  $\pbe\Renv\Penv\pred$ instantiates all subterms of $\pred$ and $\Penv$
  \label{lem:subterm-pbe}
\end{lemma}

\begin{theorem}[\textbf{Soundness \& Completeness}]
  $\pbe{\Renv}{\Penv}{\pred}$ iff $\validstar{\Renv}{\Penv}{\pred}$.
\end{theorem}
\begin{proof}
  The forward direction follows directly from the fact that
  $\pbe\Renv\Penv\pred$ performs only finitely many measure installations. For
  the reverse direction:

  \(\Renv, \Penv \vdash_* p\) means there exists an \(n\) such that
  \(\Renv, \Penv \vdash_n p\). We proceed by induction on \(n\):

  In the case where \(n=0\), we know that $\Renv, \Penv \vdash_0 p$ iff
  $\Penv \vdash p$, but in that case $\pbe\Renv\Penv\pred$,
  no matter what predicates we add to $G$.

  In the inductive case, given that
  $\forall \pred, \Renv,\Penv\vdash_n \pred \implies \pbe\Renv\Penv\pred$,
  we want to show that
  $\forall \pred, \Renv,\Penv\vdash_{n+1} \pred \implies \pbe\Renv\Penv\pred$.
  That is, there exists an expression $e$ and
  $\lambda \bar x. q \Rightarrow e \in \Renv(m)$, such that $\Penv \vdash q$ and
  $\Renv,\Penv; m  \; \bar e = e [\bar e/\bar x] \vdash_{n} \pred \implies
  \pbe\Renv\Penv\pred$

  If $\Renv,\Penv\vdash_{n}\pred$, then we're done, by the induction
  hypothesis.  Otherwise, if $\Renv, \Penv \not\vdash_{n}\pred$, but $\Renv,\Penv
  ;m\;\bar e = e [\bar e/\bar x]\vdash_{n}\pred$, then there exist
  $p_1,\ldots,p_n$, such that
  \(p_1, \ldots, p_n \not\vdash  t=t'\) but \(m\;\bar e = e[\bar e /\bar x], p_1, \ldots, p_n \vdash  t=t'\)

  By \ref{cor:subterm-exist},
  we know that $m \;\bar e$ is a subterm of $p$ or $G$,
  and then \ref{lem:subterm-pbe} gives us FIXME:
  $\pbe{\Renv;\pred}{\Penv}{m  \; \bar e = e [\bar e/\bar x]}$
  and by applying the induction hypothesis to our premise, we have
  $\pbe\Renv{\Penv}{m  \; \bar e = e [\bar e/\bar x] \rightarrow p}$
  so we have \(\pbe\Renv\Penv\pred\)
\end{proof}
